\chapter{Laryngitis}
\index{Laryngitis}

\section*{Definition and Overview}
Laryngitis is the inflammation of the larynx (voice box), which lies at the top of the airway to the lungs. It is characterised by swelling of the vocal cords, leading to hoarseness, voice loss, or throat discomfort. Laryngitis can be acute, typically lasting under three weeks and resolving spontaneously, or chronic, persisting for more than three weeks and often related to persistent irritants or underlying health conditions. The condition is a common reason for temporary vocal impairment and can significantly affect individuals whose occupations rely heavily on verbal communication.

\section*{Epidemiology}
Acute laryngitis is extremely common worldwide, frequently occurring as a component of viral upper respiratory tract infections. It is highly prevalent during winter months and affects all age groups, with young children being particularly susceptible to subglottic narrowing. Chronic laryngitis is more common in adults, particularly individuals exposed to vocal strain (such as teachers, singers, and actors) or chronic environmental irritants. Statistics suggest that acute respiratory infections remain one of the most common reasons for primary care visits in many countries, with laryngitis representing a significant proportion of self-limiting cases.

\section*{Aetiology and Pathophysiology}
Inflammation of the larynx leads to hyperaemia and oedema of the vocal fold mucosa. This swelling prevents the normal apposition and vibration of the vocal cords during phonation, altering voice pitch and quality.
\begin{itemize}
    \item \textbf{Viral infections:} the most common cause of acute laryngitis, including rhinovirus, influenza, parainfluenza, adenovirus, and coronavirus.
    \item \textbf{Vocal misuse or strain:} screaming, cheering, or prolonged speaking, particularly in loud environments, can cause acute mechanical trauma to the vocal cords.
    \item \textbf{Bacterial infections:} less common causes, including \textit{Streptococcus pneumoniae}, \textit{Haemophilus influenzae}, and \textit{Corynebacterium diphtheriae} (now rare due to immunisation).
    \item \textbf{Gastrooesophageal reflux disease (GORD):} chronic reflux of acidic stomach contents (reflux laryngitis) causes chemical irritation of the posterior larynx.
    \item \textbf{Inhaled irritants:} chronic exposure to tobacco smoke, chemical fumes, dust, allergens, or polluted air.
    \item \textbf{Corticosteroid inhalers:} improper technique when using inhaled corticosteroids for asthma can cause local irritation or secondary candidiasis.
    \item \textbf{Alcohol consumption:} high intake can cause direct irritation of the laryngeal mucosal lining.
\end{itemize}

\section*{Clinical Features}
The presentation varies from mild hoarseness to complete aphonia, often accompanied by local discomfort.
\begin{itemize}
    \item Hoarseness or a rasping, breathy voice quality, which may worsen as the day progresses
    \item Aphonia, or the complete loss of voice
    \item Sore, dry, or ticklish throat, causing a frequent tickly cough
    \item Dysphagia (difficulty swallowing) or a constant feeling of a lump in the throat
    \item Low-grade fever, rhinorrhea, and other symptoms of a viral upper respiratory tract infection
    \item Mild dyspnoea or stridor, particularly in young children, due to airway narrowing (croup)
\end{itemize}

\section*{Differential Diagnosis}
It is crucial to distinguish simple laryngitis from serious structural or neoplastic processes.
\begin{itemize}
    \item \textbf{Laryngeal neoplasia:} squamous cell carcinoma of the larynx must be excluded in any patient with hoarseness lasting more than three to four weeks, especially in smokers.
    \item \textbf{Vocal cord nodules, polyps, or cysts:} benign lesions caused by chronic vocal abuse.
    \item \textbf{Vocal cord paralysis:} neurological dysfunction affecting the recurrent laryngeal nerve, often secondary to thoracic or thyroid surgery, or neoplastic invasion.
    \item \textbf{Croup (laryngotracheobronchitis):} acute viral narrowing of the subglottic airway in young children.
    \item \textbf{Epiglottitis:} a life-threatening bacterial infection of the epiglottis causing rapid airway obstruction.
    \item \textbf{Systemic inflammatory diseases:} conditions like sarcoidosis or rheumatoid arthritis affecting the cricoarytenoid joints.
\end{itemize}

\section*{When to Seek Medical Care}
Most cases of acute laryngitis are self-limiting. However, prolonged hoarseness or signs of airway obstruction warrant prompt medical evaluation. A medical consultation is recommended if hoarseness lasts for more than three weeks, or if there is severe pain, difficulty swallowing, or coughing up blood.

\textbf{Contact your local emergency services for:}
\begin{itemize}
    \item Difficulty breathing or swallowing, or an audible high-pitched squeaking noise when breathing in (stridor)
    \item Inability to swallow saliva, leading to drooling (particularly in children)
    \item Severe chest pain or blue discolouration of the lips and face (cyanosis)
    \item High fever combined with rapid, shallow breathing
\end{itemize}

\section*{Diagnosis and Investigations}
Diagnosis is primarily clinical, based on history and voice characteristics.
\begin{itemize}
    \item \textbf{Clinical history and examination:} assessment of voice quality, duration of symptoms, and history of vocal strain or upper respiratory infection.
    \item \textbf{Laryngoscopy:} direct or indirect visualisation using a flexible nasendoscope to inspect vocal cord mobility and mucosal changes, indicated if hoarseness persists beyond three to four weeks.
    \item \textbf{Biopsy:} performed during laryngoscopy if suspicious vocal cord lesions or masses are identified.
    \item \textbf{Swabs or cultures:} rarely indicated, unless bacterial infection or fungal overgrowth (candidiasis) is suspected.
\end{itemize}

\section*{Management}
Treatment focus is on voice rest and mitigating underlying irritants.
\begin{itemize}
    \item \textbf{Voice rest:} avoiding speaking, whispering (which causes more vocal strain than quiet speaking), singing, or throat clearing.
    \item \textbf{Hydration and humidification:} drinking plenty of water and using steam inhalation to soothe irritated vocal folds.
    \item \textbf{Reflux management:} prescribing proton pump inhibitors or H2-receptor antagonists if GORD is the primary cause.
    \item \textbf{Avoidance of irritants:} complete cessation of smoking and avoiding exposure to secondhand smoke and chemical fumes.
    \item \textbf{Analgesics:} paracetamol or ibuprofen to relieve throat pain and discomfort.
    \item \textbf{Antibiotics:} not recommended for acute viral laryngitis; reserved for confirmed bacterial infections.
\end{itemize}

\section*{Complications}
\begin{itemize}
    \item Progression to chronic laryngitis due to ongoing irritation or persistent vocal strain
    \item Development of vocal cord nodules, polyps, or contact ulcers from chronic misuse
    \item Airway obstruction in children (croup) or immunocompromised patients
    \item Spread of infection to adjacent deep structures of the neck
    \item Permanent vocal changes or scarring if severe inflammatory damage goes untreated
\end{itemize}

\section*{Prognosis and Outcomes}
Acute laryngitis carries an excellent prognosis, typically resolving within one to two weeks with conservative management. Chronic laryngitis has a variable prognosis depending on the success of identifying and addressing the underlying cause (e.g., reflux control, smoking cessation, or speech therapy). Most patients regain their normal voice, though long-standing vocal cord structural changes may require specialised intervention or surgery.

\section*{Prevention and Self-Care}
\begin{itemize}
    \item Do not smoke, and avoid exposure to passive smoking.
    \item Limit alcohol and caffeine intake to prevent dehydration of the vocal tract.
    \item Stay well-hydrated by drinking six to eight glasses of water daily.
    \item Avoid vocal strain by using microphones where appropriate and avoiding yelling.
    \item Practice good hand hygiene to prevent viral upper respiratory tract infections.
    \item Use a humidifier in dry indoor environments and perform steam inhalations.
\end{itemize}

\section*{Sources and Further Reading}
The following reputable organisations provide further information on laryngitis, voice care, and treatment options.
\begin{itemize}
    \item NHS. \textit{Laryngitis causes, symptoms, and self-care guidelines.} https://www.nhs.uk/conditions/laryngitis/
    \item Mayo Clinic. \textit{Clinical overview of laryngitis diagnosis and treatment.} https://www.mayoclinic.org/diseases-conditions/laryngitis/
    \item Healthdirect Australia. \textit{Laryngitis and voice loss health advice.} https://www.healthdirect.gov.au/laryngitis
    \item Better Health Channel. \textit{Information on laryngitis and voice care.} https://www.betterhealth.vic.gov.au/health/conditionsandtreatments/laryngitis
    \item MSD Manual. \textit{Professional reference page on acute and chronic laryngitis.} https://www.msdmanuals.com/
\end{itemize}
